\documentclass[a4paper,11pt]{article}

% ==============================================================================
% CHARGEMENT DU STYLE DU TEMPLATE (template.sty)
% ==============================================================================
\usepackage{template}

% ==============================================================================
% INFORMATIONS DU RAPPORT (À PERSONNALISER POUR CHAQUE RENDU)
% ==============================================================================
\matiere{Nom de la Matière}                           % Affiché sur la couverture et dans l'en-tête (darkred)
\titre{Titre du Rapport}                               % Titre principal du rapport
\soustitre{Sous-titre ou intitulé du projet}           % Sous-titre (laisser vide {} si aucun)
\auteur{Prénom Nom}                                    % Votre nom (ou plusieurs séparés par \and ou une virgule)
\formation{Master Informatique}                        % Intitulé de votre formation
\etablissement{Université}                             % Nom de l'établissement / Université
\annee{2026 -- 2027}                                   % Année universitaire ou date de rendu

% ==============================================================================
% DÉBUT DU DOCUMENT
% ==============================================================================
\begin{document}

% Génération automatique de la page de garde
\makecoverpage

% ==============================================================================
% BLOCS PRÉFAITS ET EXEMPLES (DÉCOMMENTER LES BLOCS SOUHAITÉS)
% ==============================================================================

% --- TABLE DES MATIÈRES ---
% \tableofcontents
% \newpage

% --- STRUCTURE DU DOCUMENT (SECTIONS & SOUS-SECTIONS) ---
% \section{Titre de la section}
% Contenu du paragraphe...
%
% \subsection{Titre de la sous-section}
% Contenu du sous-paragraphe...
%
% \subsubsection{Titre de la sous-sous-section}
% Détails techniques...

% --- BOÎTES DE MESSAGE PRÉFAITES ---
%
% 1. Infobox (Information générale - Bleu)
% \begin{infobox}[Titre de l'information (optionnel)]
%   Contenu de l'infobox...
% \end{infobox}
%
% 2. Warningbox (Avertissement / Attention - Ambre)
% \begin{warningbox}[Attention (optionnel)]
%   Contenu de l'avertissement...
% \end{warningbox}
%
% 3. Notebox (Note / Remarque - Darkred, couleur du thème)
% \begin{notebox}[Remarque (optionnel)]
%   Contenu de la note...
% \end{notebox}
%
% 4. Tipbox (Astuce / Conseil - Vert)
% \begin{tipbox}[Astuce (optionnel)]
%   Contenu du conseil ou de la bonne pratique...
% \end{tipbox}
%
% 5. Dangerbox (Danger / Erreur critique - Rouge vif)
% \begin{dangerbox}[Danger critique (optionnel)]
%   Contenu du message critique...
% \end{dangerbox}

% --- LISTES À PUCES ET NUMÉROTÉES ---
% \begin{itemize}
%   \item Premier élément
%   \item Deuxième élément
% \end{itemize}
%
% \begin{enumerate}
%   \item Première étape
%   \item Deuxième étape
% \end{enumerate}

% --- INSERTION D'UNE FIGURE ---
% \begin{figure}[htbp]
%   \centering
%   \includegraphics[width=0.75\textwidth]{nom_de_l_image.png}
%   \caption{Description de la figure}
%   \label{fig:mon_image}
% \end{figure}

% --- INSERTION D'UN TABLEAU ---
% \begin{table}[htbp]
%   \centering
%   \caption{Description du tableau}
%   \label{tab:mon_tableau}
%   \begin{tabular}{llr}
%     \toprule
%     \textbf{Colonne 1} & \textbf{Colonne 2} & \textbf{Valeur} \\
%     \midrule
%     Donnée A & Détail A & 12.5 \\
%     Donnée B & Détail B & 45.0 \\
%     \bottomrule
%   \end{tabular}
% \end{table}

% --- EXTRAIT DE CODE SOURCE ---
% \begin{lstlisting}[language=Python, caption={Exemple de script Python}]
% def exemple():
%     # Commentaire
%     return True
% \end{lstlisting}

% --- FORMULES MATHÉMATIQUES ---
% Formule en ligne : $E = mc^2$
% Formule centrée numérotée :
% \begin{equation}
%   f(x) = \int_{0}^{+\infty} e^{-t} t^{x-1} \, \mathrm{d}t
% \end{equation}

\end{document}
